% =============================================================================
% Q2 -- BRIDGE DRAFT (2026-08-23), between phrase-blocks and full prose.
% Not final chapter text -- connected clauses/fragments with numbers, figures, and
% tables woven in inline, so nothing load-bearing gets lost. Write your own sentences
% from this directly; do not treat this file as compilable prose.
% Ordering: headline-first, per the agreed Q2 story (weak result folds into the real
% headline as a pivot, not a separate opening section).
% Source tables/figures already built:
%   Table~\ref{tab:results-q2}            -- old CI table (4/6survey, EN/XGB/GNNWR)
%   Table~\ref{tab:results-q2-ablation}   -- CanopyCover ablation grid (NEW)
%   Table~\ref{tab:results-q2-reliability}-- Q1/Q2 variable-reliability cross-check (NEW)
%   Figure~\ref{fig:q2-example-plot}      -- q2_example_plot_curve.png
%   Figure~\ref{fig:q2-example-plot-typical} -- q2_example_plot_curve_typical_appendix.png (appendix)
%   Figure~\ref{fig:q2-stability-map}     -- q2_topex_vs_canopycover_stability.png
% =============================================================================

--- 0. OPENING / PIVOT ---
Q1: global models (one formula, whole forest) -> real spatial clustering left in errors.
Q2 test: does letting the model vary by location (GNNWR) pick up what global misses?

Point estimate: GNNWR > EN, XGBoost, every set (Table~\ref{tab:results-q2}) --
Set4: GNNWR 0.294 [0.236, 0.347] vs. EN 0.240 [0.192, 0.286] vs. XGBoost 0.250 [0.201, 0.295].
But: CIs overlap every set -- not statistically distinguishable.
Also modest in absolute terms -- best point estimate (Set2, 0.320) still leaves ~2/3 of
Delta y_max unexplained.
Pivot needed here: tested WHY GNNWR wins even by that small, uncertain margin -- answer
overturns the reason for trying GNNWR in the first place.

--- 1. THE HEADLINE -- CanopyCover explains the whole edge ---
Claim: GNNWR's entire advantage over the global models = one variable (CanopyCover), not
a general ability to read environmental relationships by place.

Test: ablation -- remove CanopyCover, refit all three (Table~\ref{tab:results-q2-ablation}).
  full Set4:       EN 0.240 / XGB 0.250 / GNNWR 0.294 (0.277+/-0.068 fold-mean)
  no CanopyCover:  EN 0.042 / XGB 0.061 / GNNWR 0.049 (0.026+/-0.067 fold-mean) -- all
                   three collapse to the same near-zero floor
  CanopyCover ONLY, nothing else: GNNWR 0.249 (0.235+/-0.030 fold-mean) -- recovers
                   ~82% of the full gap: (0.249-0.049)/(0.294-0.049)

Ruled out: circularity -- raw correlation CanopyCover vs. target = 0.468 (r^2~0.22),
moderate, not near 1 -- but still means nothing else in the 19-variable set carries
independent weight.

So-what: the hypothesis behind testing GNNWR (spatial modelling reveals hidden
environmental drivers) is wrong -- GNNWR just re-found the one variable every method
already had, expressed slightly differently by place.

--- 2. WHY SO LITTLE ELSE COMES THROUGH -- the target is noisier than it looks ---
Claim: Delta y_max isn't a clean, isolated ceiling measurement -- it's mixed with
growth-rate mismatch.

Mechanism: curve-fit holds each plot's growth rate (k) fixed to its yield class, only
fits the ceiling -- if real rate != assumed rate, the mismatch doesn't disappear, it
leaks into the fitted ceiling instead.

Worked example -- Figure~\ref{fig:q2-example-plot}, plot 309770 (mature stand, observed
ages 65-81): own fitted ceiling 43.0m vs. yield-class benchmark 32.8m -> Delta y_max =
+10.2m; XGBoost predicts +12.4m for this plot; shaded region on the trajectory panel =
the real, measured gap over the observed ages (not just the extrapolated ceiling
difference).
Worth a phrase: CanopyCover's SHAP contribution is NEGATIVE for this specific plot even
though it's positive on average across the population -- direction isn't fixed per plot.

Appendix companion -- Figure~\ref{fig:q2-example-plot-typical}, plot 166972 (picked by
median |residual|, not cherry-picked): actual +11.2m, model predicts +5.3m, roughly half
explained, no single dominant SHAP driver -- discloses that Figure 1's plot is unusually
well-explained, not typical.

So-what: explains finding 1 mechanistically -- a noisy target leaves little clean signal
for anything except the variable most directly tied to real, measured structure
(CanopyCover).

--- 3. CROSS-CHECK WITH Q1 -- same variables, same behaviour, two different tests ---
Claim: terrain-variable reliability pattern matches Q1's -- found via a genuinely
different method each time (Table~\ref{tab:results-q2-reliability}).

Q2's test: GNNWR's own local coefficient, % of compartments with the same sign (n=231) --
CanopyCover 100.0%, slope_degrees 97.0%, chelsa_bio12_precip_mm 87.4%, topex 71.9%.

Q1's test: does a global Elastic Net model agree with a linear mixed-effects model
(LMM) on a variable's coefficient? slope_degrees & precip: EN/LMM agree (both positive /
both negative respectively) -> reliable. topex: EN~=0, LMM=+1.32 -> disagree -- BUT a
mechanism check (within-compartment correlation vs. raw correlation) shows this
disagreement is because topex carries REAL local signal (within-corr 0.214 > raw-corr
0.080) that a single straight line struggles to capture, not because it's noise.

Caveat, must state plainly: these are two different measurements (global-model agreement
vs. local-coefficient sign stability) that happen to point the same way -- not the same
test literally repeated. The agreement is more informative because of that, not less.

Visual -- Figure~\ref{fig:q2-stability-map}: CanopyCover map (uniformly stable, one
colour direction) vs. topex map (mixed sign, patchy) side by side, same colour scale.

So-what: strengthens confidence this isn't a one-chapter artifact -- the same variables
behave consistently regardless of method or target.

--- 4. LOOSE ENDS -- one phrase-chain each, not developed into their own paragraphs ---
Terrain non-linear: windward_topex & slope_degrees show a real, saturating (rises then
flattens) SHAP relationship to the target that XGBoost can see and GNNWR structurally
can't fully use (GNNWR is local-*linear* -- one straight line per place) -- a small
effect overall, since CanopyCover already carries ~80-90% of the model's total
explanatory power.

Outliers: under 1% of the 4-survey cohort (266/56,526), Tukey-fence flagged -- same
rate/ceiling-mixing mechanism as finding 2, at an extreme scale -- not a separate
headline; full diagnostics (71.6% artifact-driven / 28.4% genuinely fast growth split,
exposed vs. sheltered terrain) preserved as an appendix-candidate block, not developed
in the main text given how small a share of the data this is.

6-survey: run as a planned sensitivity check, not the primary analysis -- R^2~0 for
Elastic Net, XGBoost, and GNNWR alike, consistent with far fewer distinct compartments
(47 vs. 231) -- not investigated further.

--- 5. CLOSING -- so what for the dissertation ---
Genuine negative result, not a failed test -- tells you the LIMIT of this approach
(a target that blends growth-rate and ceiling noise together, plus a spatially-varying
model that can only draw a straight line per place), not just that it didn't work.
Contrast phrase, to fill in once Q3 is drafted: Q2's negative result vs. Q3's positive
one -- worth stating why Q3 delivered practical value and Q2 delivered a boundary
condition instead.
